
\question[3]
Erkläre, warum Programmcode und Textdokumente ausschließlich verlustfrei 
komprimiert werden dürfen, während bei Videos oder Musikdateien verlustbehaftete 
Kompression üblich ist.

\begin{solutionbox}{5cm}
Bei Programmcode und Textdokumenten führt schon die Änderung eines einzelnen Zeichens
(z.\,B. eines Semikolons im Quellcode oder eines Buchstabens in einem Vertragstext) 
zu einer fehlerhaften bzw. unbrauchbaren Datei -- der Code lässt sich dann eventuell
nicht mehr ausführen oder der Text verändert seine Bedeutung. Diese Daten müssen deshalb
bit-genau erhalten bleiben.
Bei Videos und Musik hingegen nimmt der Mensch kleine Abweichungen 
(leichte Farbänderungen, kaum hörbare Frequenzanteile) meist nicht oder kaum wahr.
Deshalb kann hier zugunsten einer wesentlich kleineren Dateigröße ein geringer, 
kaum bemerkbarer Qualitätsverlust in Kauf genommen werden.
\end{solutionbox}